% !TeX root = ../papers/nc-p3.tex
\section{接する二成分と三直線を含む被約三次判別式の全持上げ}
\label{sec:marked-reduced-cubic-commutators}
\begin{theorem}\label{thm:marked-reduced-cubic-commutators}
標数零の代数閉体上、判別式が任意の被約な平面三次曲線である順序付き跡零三行列 $C$ の
全三交換子逆像が以下と先の各判別式型の手順で得られる。直線成分に階数一補正を持つ解は各直線の
圧縮対と四変数の階数層から復元する。補正を一つも持たない解について、直線とそれに接する
滑らかな二次曲線、三角形、共点三直線の場合に残る補助表示は、
それぞれ局所自由で多重次数 $(1,2),(1,1,1),(1,1,1)$ の表示だけである。
目標も同じ型でなければこの部分は空である。この型なら三角形で二つ、それ以外で一つの
補助候補があり、各候補の正確な逆像は線型に得る対称行列の階数3,2,1に従って
二点、空、滑らかなアフィン二次曲面である。

直線と接する二次曲線で局所自由な $(1,2)$ 型の一般の目標の全逆像は、
二つの孤立解と、各々 $\mathbf A^1$ に同型な二つのアフィン二次曲線である。
\end{theorem}
\begin{proof}[補助表示の線束と成分次数]
線束の貼合せと双一次関係で決定式の行空間を得る方法は
定理\ref{thm:marked-singular-integral-commutators}と同じである。
直線で左核が一定なら補題\ref{lem:marked-squarefree-symmetrizer}が階数一補正を与える。
全成分が直線で左核がいずれも動くとする。左核を与える余核層の正規化上の線束の
各次数は1以上、その和は3以下である。後者を確かめるには各特異点の完備局所環 $R$ で
加群の分数イデアルを、正規化を生成し単位1を含む $I$ に規格化する。
正規化を生成すれば無限体上で全ての枝で単元な一般元を選べる。
線束として貼り合わせた場合とのEuler標数の差から、正規化の次数和は
$3-\sum\operatorname{length}(I/R)$ となる。従って等号で全て局所自由、
三つの次数は全て1である。余因子の左右核積から右側も同じ次数である。

接する場合は $E:x(xz-y^2)=0$、成分を $(0:s:t),(s^2:st:t^2)$ とする。
接点で $R$ は両枝の値と一次ジェットが一致する部分環で、$\bar R/R$ は
共通ジェット $k[r]/(r^2)$ 上の長さ二の巡回加群である。従って $R\subset I\subset\bar R$
で中間にあるのは値だけを合わせる節点環一つである。
多重次数の候補は、前節の二交点を順に分離した候補と同じである。
実際導手に沿う消失は今度は相異なる二点での消失または単一点での二次消失であり、
$\mathbf P^1$ 上の切断数と二つのジェット評価階数は同じである。
余因子の直線上の共通零の次数は三局所型で順に $0,1,2$。例えば中間の節点環は
二枝の値を合わせるので二つの生成元で表示され、平面局所座標 $u,v$ で
枝を $v=0,v=u^2$ とすれば、基底 $(1,1),(0,u)$ に対する $v$ の作用は
$\left(\begin{smallmatrix}0&0\\u&u^2\end{smallmatrix}\right)$ である。
直線 $v=0$ 上の非単位ブロックの成分の共通零は一次となる。
完全正規化の場合は二枝の対角表示で次数2を得る。

左核が動き右核が直線上で定点の補助解は、自由な $(0,3)$ と中間の $(0,2)$
の場合だけである。順に代表は
\[
 \begin{pmatrix}y&-x&0\\z&-y&0\\0&z-e y&-x\end{pmatrix},\qquad
 \begin{pmatrix}y&-x&0\\0&y&-x\\z&0&-x\end{pmatrix}.
\]
前者の右核は二次曲線上で $(s^3:s^2t:t^3-e st^2)$ であり、後者では $(s^2:st:t^2)$、
直線上はいずれも $e_2$ である。前者のジェット倍率 $e$ は $s/t$ の倍率で $0,1$ にできる。
いずれも六線型解の全行列が $ke_2$ を保ち、非零判別式を生じない。
以上で補正のない解は主張の三型に限られる。
\end{proof}

\begin{proof}[三種類の余因子公式]
接する二成分の自由な $(1,2)$ の右核は、二次曲線上 $(s^2:st:t^2)$、
直線上 $(0:s:t+e s)$ と書ける。$e\ne0$ であり、$e=0$ は平面線束そのものになる。
倍率で $e=1$ とする。以下に補助行列 $A$ と目標の参照行列 $M$ を並べる。
\[
 A=\begin{pmatrix}y&-x&0\\z&0&-x\\0&y+z&-x-y\end{pmatrix},
 \qquad M=\begin{pmatrix}x+y&-x&0\\-x+z&2x&-x\\0&-x+z&x-y\end{pmatrix}.
\]
$M$ の直線上と二次曲線上の右核は $(0:s:t)$ と $(s^2:s(s+t):(s+t)^2)$ であり、
直線の正規化と二次曲線の正規化でのジェットの差は $-2$ となる。

三角形では各頂点で左右二枝の値を貼り合わせる。各成分のスカラーで二頂点の倍率を
規格化すれば、残る一つの倍率が平面線束に対する相対値を分類する。代表は
\[
 A=\begin{pmatrix}x&\lambda z&0\\0&y&x\\y&0&z\end{pmatrix},
 \qquad M=\begin{pmatrix}x&-z&0\\0&\lambda y&-x\\y&0&-\lambda z\end{pmatrix}.
\]
ここで $\lambda\ne0,-1$。三辺で右核の非零二座標を取り、頂点で巡回して比を掛けると
$A$ の相対倍率は $-\lambda$（巡回を逆にすればその逆数）である。
$M$ の行列式の二単項積の係数比は $-1/\lambda^2$ であり、
対角行列による左右の基底変更でこの値を $\lambda$ とする $A$ の形にできる。
従って二候補となる。$\lambda=1$ では $\det M=0$ で補助解はない。

共点三直線を $(r:0:1),(0:r:1),(r:r:1)$ で取る。平面の局所関数は三枝の値が一致し、
第三枝の一次係数が前二枝の和である。二次以上では各枝の係数を独立に指定できる。
自由な多重次数 $(1,1,1)$ の線束は値を一致させた後、第三の一次係数を前二つの和に
$e$ 倍の共通の値を加えたものにする。$e\ne0$ を倍率で1にしてよい。
三右核の埋込みは $(r:0:1),(0:r:1),(r:r:1+r)$ であり、代表は
\[
 A=\begin{pmatrix}y&-x&0\\z&x&-x\\0&x+z&-y\end{pmatrix},
 \qquad M=\begin{pmatrix}y&-x&0\\x+z&-x&-x\\0&-x+y+z&-y\end{pmatrix}.
\]
$M$ の三右核は $(r:0:1+r),(0:r:1+r),(r:r:1)$ なので相対ジェットは $-2$ である。
二つの加法型で倍率を戻せば $e\mapsto-2e$ となり候補は一意である。

三表示とも六空間を $Z$ によって線型座標化すると
$C_{Q(Z)}=\Psi(\operatorname{adj}Z)M$ を得る。再現のため、対角二つと六つの
非対角による跡零の基底を、行優先の行列単位からスカラー部分を引き $E_{22}$ を除いて作る。
$Q_0,Q_1,Q_2$ のこの順に並べた二十四変数で $\operatorname{im}\Gamma_Q\subset W$ を
簡約行階段形にし、自由列の順に変数を置く。前節の $Z=(P,T,U,A',B,D)$ の記号で
（$A'$ は $01$ 成分）、その六自由変数は三型の順に
\[
 (U,-T+D,-D,A'-B,B,P),\quad
 (U,-D,(\lambda-1)A',T,-B,\lambda P),\quad
 (U,-T,-D,P,B,A').
\]
三角形ではまず不定元 $\lambda$ で計算する。得る $Q(Z)$ は多項式で、
$\lambda=1$ でも線型階数6である。他は自由列指定から同型が分かる。
左倍率の六単項式抽出には順に $(1,1,0),(1,1,1),(1,2,1)$ で評価する。
六列の小行列式は順に $-1,-\lambda^5,-1$。各係数の恒等式と特殊化の計算は
\path{computations/verify_reduced_cubic_commutators.py} にある。
$M$ の具体的な右核像は二次曲線または複数の異なる直線を含み全点を固定する
射影変換は恒等なので、従前の固定目標の左右比較と余因子逆像がそのまま使える。
\end{proof}

\begin{proof}[補正を持つ解と接する一般ファイバー]
直線ごとに前節の $K,kw$ を求め、前節の四変数の階数層を用いる。
三直線で複数の補正がある場合は和集合を取り、または直線を順序付けて最初の
補正のみを採用する。核が両方定点の場合の二行二列商鉛筆が比例する場合もここでは起こり得る。
そのとき可逆スライスで割った他の作用がスカラーなので、$K/kr$ は
射影直線の全点を動かす。一般スライスが階数2であることは被約性から従う。
これで交点の階数が零になる場合も含め全候補を得る。

接する二成分の一般ファイバーの型を調べる。前節の半単純な一旗の解法で
$G=hC_0$ とする。一般座標で $\alpha_1=(1,0),\alpha_2=(0,1)$、
$u_1=(0,m)^t,u_2=(n,0)^t$ としてよい。実際二行と二列が各々独立で、
$(\alpha_i u_j)$ の対称部分が非退化なら、その二つの等方直線を添字の基底に取り、
二行を平面の双対基底にできる。$J=\left(\begin{smallmatrix}j&k\\l&-j\end{smallmatrix}\right)$
と置くと $G_{11}=J+\operatorname{diag}(-n,m)$。直線での残余二次曲線の式は
符号・非零倍率と直線座標の変更まで
\[
              -mk s^2+j(m-n)st-nl t^2 .
\]
実際 $(s,t)\operatorname{adj}(G_{11})(nt,ms)^t$ を展開すればよい。
一方解のアフィン平面の無限遠で $c=(D_1)_{10}$、$f=(D_1)_{11}$ とする。
線型条件により $(D_1)_{01}=-kc/l+j(n-m)f/(nl)$、
$(D_2)_{10}=j(n-m)c/(mk)-lf/k$、$(D_2)_{01}=f$ となる。
従って $[D_1,D_2]_{00}=0$ の無限遠式は非零倍率まで
\[
                   nk c^2-j(n-m)cf+ml f^2 .
\]
二式の判別式は一致する。ゆえに接する場合は無限遠点が一つである。
一般条件の実現を $n=2,m=1,j=4,k=1,l=2$ と
$G=\left(\begin{smallmatrix}1&2&-2\\-2&2&1\\-3&2&-3\end{smallmatrix}\right)$ で見る。
残余二次曲線は $-30x^2-15xy-2xz-(y+2z)^2$ で滑らか、$h=1$、
解の平面の式は $(g-v)^2-10g-4v+11=0$ で滑らかである。
被約な接する成分対は一軌道で、自由な $(1,2)$ の目標は貼合せ、右基底、可逆な跡零左倍率から
既約である。よってこの非空な開条件と余因子の階数3が同時に一般点で成立し、
二符号それぞれで射影二次曲線から一点を除いたものを得る。
\end{proof}
